\documentclass{article}
\usepackage[dblblindworkshop]{neurips_2026}
\workshoptitle{Geometric Distributional Deep Learning: Bridging Optimal Transport, Learning and Structured Data}
\usepackage[utf8]{inputenc}
\usepackage[T1]{fontenc}
\usepackage{amsmath,amssymb,amsthm}
\usepackage{hyperref}
\hypersetup{hidelinks,pdftitle={},pdfauthor={},pdfsubject={},pdfkeywords={}}
\usepackage{url}

\title{Anonymous supplement: rational VP Heun certificate}
\author{Anonymous Author(s)}

\begin{document}
\maketitle

This supplement records the integer comparison used in the workshop
paper. It contains no author names, emails, or public repository URLs.

The four-step VP Heun product is the following element $P(\pi,\sqrt{2})$
of $\mathbb{Q}[\pi,\sqrt{2}]$, obtained by expanding the four Heun
factors with $h=1/4$ and reducing $\sqrt{2}^2=2$:
\begin{equation}
\begin{aligned}
P(\pi,\sqrt{2})
&= 1+\frac{3}{40}\pi+\frac{15}{164}\pi\sqrt{2}
+\frac{78579}{10758400}\pi^2+\frac{45}{5248}\pi^2\sqrt{2}\\
&\quad+\frac{11421}{17213440}\pi^3
+\frac{4671}{17213440}\pi^3\sqrt{2}\\
&\quad+\frac{5103}{275415040}\pi^4
+\frac{567}{55083008}\pi^4\sqrt{2}\\
&\quad+\frac{243}{2203320320}\pi^5
+\frac{729}{2203320320}\pi^5\sqrt{2}
+\frac{729}{176265625600}\pi^6.
\end{aligned}
\end{equation}
Each of the twelve coefficients is a ratio of nonnegative integers,
hence nonnegative, and the constant term is $1$. Substituting
$\pi<355/113$ and $\sqrt{2}<99/70$ is therefore valid and yields
\[
P\Bigl(\frac{355}{113},\frac{99}{70}\Bigr)
=\frac{192113353671412470139303}{102753427879939708813312}
<\frac{187}{100},
\]
equivalently
\[
19211335367141247013930300<19214891013548725548089344.
\]
Hence $r_{\mathrm{VP}}<187/100$ and $W_2^{\mathrm{VP}}>13/100$.
The bound $\sqrt{2}<99/70$ is $99^2-2\cdot 70^2=1$. The bound
$\pi<355/113$ follows from Machin's formula with Leibniz remainders
(seven Taylor terms for $\arctan(1/5)$, four for $\arctan(1/239)$).
An independent \texttt{mpmath}~$1.3.0$ interval enclosure is a
cross-check, not the proof.

\end{document}
